
\begin{filecontents}{refs.bib}
    @article{2009,
    title={2009 10th International Conference on Document Analysis and Recognition Hybrid Page Layout Analysis via Tab-Stop Detection}
    }
    
    @article{Ashwin2002A,
    title={A font and size-independent OCR system for printed Kannada documents using support vector machines},
    author={T. Ashwin and P. Sastry},
    journal={Sadhana},
    year={2002},
    volume={27},
    pages={35-58},
    doi={10.1007/bf02703311}
    }
    
    @article{Bhandari2026A,
    title={A Framework and Dataset for Contextual Post-OCR Correction},
    author={Abhishek Bhandari and Gaurav Harit},
    journal={ACM Transactions on Asian and Low-Resource Language Information Processing},
    year={2026},
    volume={25},
    pages={1 - 20},
    doi={10.1145/3815575}
    }
    
    @article{Caravani2026Assessing,
    title={Assessing and comparing document OCR systems in the era of Large Language Models},
    author={Valerio Caravani and Riccardo De Cesaris and P. Merialdo},
    journal={Inf. Process. Manag.},
    year={2026},
    volume={64},
    pages={105047},
    doi={10.1016/j.ipm.2026.105047}
    }
    
    @article{Chaudhuri2002Page,
    title={Page layout analyser for multilingual Indian documents},
    author={A. R. Chaudhuri and A. Mandal and B. B. Chaudhuri},
    journal={Language Engineering Conference, 2002. Proceedings},
    year={2002},
    pages={24-32},
    doi={10.1109/lec.2002.1182287}
    }
    
    @article{Chaudhuri1997An,
    title={An OCR system to read two Indian language scripts: Bangla and Devnagari (Hindi)},
    author={B. Chaudhuri and U. Pal},
    journal={Proceedings of the Fourth International Conference on Document Analysis and Recognition},
    year={1997},
    volume={2},
    pages={1011-1015 vol.2},
    doi={10.1109/icdar.1997.620662}
    }
    
    @article{Faraz2026Designing,
    title={Designing Production-Scale OCR for India: Multilingual and Domain-Specific Systems},
    author={Ali Faraz and Raja Kolla and Ashish Kulkarni and S. Agarwal},
    journal={ArXiv},
    year={2026},
    volume={abs/2602.16430},
    doi={10.48550/arxiv.2602.16430}
    }
    
    @article{Faraz2025IndicVisionBench:,
    title={IndicVisionBench: Benchmarking Cultural and Multilingual Understanding in VLMs},
    author={Ali Faraz and Akash and Shaharukh Khan and Raja Kolla and Akshat Patidar and Suranjan Goswami and Abhinav Ravi and Chandra Khatri and Shubham Agarwal},
    journal={ArXiv},
    year={2025},
    volume={abs/2511.04727},
    doi={10.48550/arxiv.2511.04727}
    }
    
    @article{Fateh2023Enhancing,
    title={Enhancing optical character recognition: Efficient techniques for document layout analysis and text line detection},
    author={A. Fateh and M. Fateh and Vahid Abolghasemi},
    journal={Engineering Reports},
    year={2023},
    volume={6},
    doi={10.1002/eng2.12832}
    }
    
    @article{Govindaraju2004Tools,
    title={Tools for enabling digital access to multi-lingual Indic documents},
    author={V. Govindaraju and Swapnil Khedekar and Suryaprakash Kompalli and Faisal Farooq and Srirangaraj Setlur and V. Ramanaprasad},
    journal={First International Workshop on Document Image Analysis for Libraries, 2004. Proceedings.},
    year={2004},
    pages={122-133},
    doi={10.1109/dial.2004.1263244}
    }
    
    @article{Hamdi2025PILOT:,
    title={PILOT: A promptable interleaved layout-aware OCR transformer},
    author={Laziz Hamdi and Amine Tamasna and P. Boisson and Thierry Paquet},
    journal={International Journal on Document Analysis and Recognition (IJDAR)},
    year={2025},
    doi={10.1007/s10032-026-00590-w}
    }
    
    @article{Jayatilleke2025Zero-shot,
    title={Zero-shot OCR Accuracy of Low-Resourced Languages: A Comparative Analysis on Sinhala and Tamil},
    author={Nevidu Jayatilleke and Nisansa de Silva},
    year={2025},
    pages={471-480},
    doi={10.48550/arxiv.2507.18264}
    }
    
    @article{Jhala2026PragyaDoc:,
    title={PragyaDoc: A Universal Document Intelligence Framework for Multilingual Medical Document Understanding in Low-Resource Settings},
    author={Jagpal Singh Jhala},
    year={2026},
    doi={10.48550/arxiv.2608.07478}
    }
    
    @article{Kashid2024RoundTripOCR:,
    title={RoundTripOCR: A Data Generation Technique for Enhancing Post-OCR Error Correction in Low-Resource Devanagari Languages},
    author={Harshvivek Kashid and Pushpak Bhattacharyya},
    journal={ArXiv},
    year={2024},
    volume={abs/2412.15248},
    doi={10.48550/arxiv.2412.15248}
    }
    
    @article{Krishnan2014Towards,
    title={Towards a Robust OCR System for Indic Scripts},
    author={Praveen Krishnan and Naveen Sankaran and A. Singh and C. Jawahar},
    journal={2014 11th IAPR International Workshop on Document Analysis Systems},
    year={2014},
    pages={141-145},
    doi={10.1109/das.2014.74}
    }
    
    @article{Kumar2012Discrimination,
    title={Discrimination of English to other Indian languages (Kannada and Hindi) for OCR system},
    author={A. Kumar and Tushar Patnaik and V. K. Verma},
    journal={ArXiv},
    year={2012},
    volume={abs/1205.2164},
    doi={10.5121/ijcsea.2012.2214}
    }
    
    @article{Li2025Printed,
    title={Printed document layout analysis and optical character recognition system based on deep learning},
    author={Dong-Lin Li and Shih-Kai Lee and Yin-Ting Liu},
    journal={Scientific Reports},
    year={2025},
    volume={15},
    doi={10.1038/s41598-025-07439-y}
    }
    
    @article{Li2026Structure-Aware,
    title={Structure-Aware Multi-Task Learning for Robust Vision-Language OCR},
    author={Zhe-Han Li},
    journal={2026 IEEE 8th International Conference on Communications, Information System and Computer Engineering (CISCE)},
    year={2026},
    pages={638-643},
    doi={10.1109/cisce69494.2026.11504766}
    }
    
    @article{Madhani2023Bhasa-Abhijnaanam:,
    title={Bhasa-Abhijnaanam: Native-script and romanized Language Identification for 22 Indic languages},
    author={Yash Madhani and Mitesh M. Khapra and Anoop Kunchukuttan},
    year={2023},
    pages={816-826},
    doi={10.48550/arxiv.2305.15814}
    }
    
    @article{Maheshwari2022A,
    title={A Benchmark and Dataset for Post-OCR text correction in Sanskrit},
    author={Ayush Maheshwari and Amrith Krishna and Ganesh Ramakrishnan},
    journal={ArXiv},
    year={2022},
    volume={abs/2211.07980},
    doi={10.48550/arxiv.2211.07980}
    }
    
    @article{Malik2026synthocr-gen:,
    title={synthocr-gen: A synthetic ocr dataset generator for low-resource languages- breaking the data barrier},
    author={Haq Nawaz Malik and K. Shafi and Tanveer Ahmad Reshi},
    journal={ArXiv},
    year={2026},
    volume={abs/2601.16113},
    doi={10.48550/arxiv.2601.16113}
    }
    
    @article{Manukonda2025byteSizedLLM@NLU,
    title={byteSizedLLM@NLU of Devanagari Script Languages 2025: Language Identification Using Customized Attention BiLSTM and XLM-RoBERTa base Embeddings},
    author={D. Manukonda and R. Kodali},
    year={2025},
    pages={248-252}
    }
    
    @article{Mathew2016Multilingual,
    title={Multilingual OCR for Indic Scripts},
    author={Minesh Mathew and A. Singh and C. Jawahar},
    journal={2016 12th IAPR Workshop on Document Analysis Systems (DAS)},
    year={2016},
    pages={186-191},
    doi={10.1109/das.2016.68}
    }
    
    @article{Mishra2014WORD,
    title={WORD RECOGNITION IN INDIC SCRIPTS},
    author={Anand Mishra and Praveen Krishnan and R. Shekhar},
    year={2014}
    }
    
    @article{Nath2025IndicDLP:,
    title={IndicDLP: A Foundational Dataset for Multi-lingual and Multi-domain Document Layout Parsing},
    author={Oikantik Nath and Sahithi Kukkala and Mitesh M. Khapra and Ravi Kiran Sarvadevabhatla},
    journal={ArXiv},
    year={2025},
    volume={abs/2512.20236},
    doi={10.1007/978-3-032-04614-7_2}
    }
    
    @article{Pal2004Indian,
    title={Indian script character recognition: a survey},
    author={U. Pal and B. Chaudhuri},
    journal={Pattern Recognit.},
    year={2004},
    volume={37},
    pages={1887-1899},
    doi={10.1016/j.patcog.2004.02.003}
    }
    
    @article{Pati2004Gabor,
    title={Gabor filters for document analysis in Indian bilingual documents},
    author={P. B. Pati and S. Raju and N. Pati and A. Ramakrishnan},
    journal={International Conference on Intelligent Sensing and Information Processing, 2004. Proceedings of},
    year={2004},
    pages={123-126},
    doi={10.1109/icisip.2004.1287637}
    }
    
    @article{Qiao2024Reading,
    title={Reading order detection in visually-rich documents with multi-modal layout-aware relation prediction},
    author={Liang Qiao and Can Li and Zhanzhan Cheng and Yunlu Xu and Yi Niu and Xi Li},
    journal={Pattern Recognit.},
    year={2024},
    volume={150},
    pages={110314},
    doi={10.1016/j.patcog.2024.110314}
    }
    
    @article{Sahare2018Multilingual,
    title={Multilingual Character Segmentation and Recognition Schemes for Indian Document Images},
    author={Parul Sahare and S. B. Dhok},
    journal={IEEE Access},
    year={2018},
    volume={6},
    pages={10603-10617},
    doi={10.1109/access.2018.2795104}
    }
    
    @article{Saluja2017A,
    title={A Framework for Document Specific Error Detection and Corrections in Indic OCR},
    author={Rohit Saluja and D. Adiga and Ganesh Ramakrishnan and Parag Chaudhuri and Mark James Carman},
    journal={2017 14th IAPR International Conference on Document Analysis and Recognition (ICDAR)},
    year={2017},
    volume={04},
    pages={25-30},
    doi={10.1109/icdar.2017.308}
    }
    
    @article{Saluja2017Error,
    title={Error Detection and Corrections in Indic OCR Using LSTMs},
    author={Rohit Saluja and D. Adiga and Parag Chaudhuri and Ganesh Ramakrishnan and Mark James Carman},
    journal={2017 14th IAPR International Conference on Document Analysis and Recognition (ICDAR)},
    year={2017},
    volume={01},
    pages={17-22},
    doi={10.1109/icdar.2017.13}
    }
    
    @article{Sarungbam2014Script,
    title={Script identification and language detection of 12 Indian languages using DWT and template matching of Frequently Occurring Character(s)},
    author={Jeelen Kumar Sarungbam and B. Kumar and Ankur Choudhary},
    journal={2014 5th International Conference - Confluence The Next Generation Information Technology Summit (Confluence)},
    year={2014},
    pages={669-674},
    doi={10.1109/confluence.2014.6949300}
    }
    
    @article{Sharma2019Categorization,
    title={Categorization of Multilingual Text on Languages of Indic Script},
    author={Dr. H. R. Sharma and Sreebha Bhaskaran},
    journal={2019 Global Conference for Advancement in Technology (GCAT)},
    year={2019},
    pages={1-4},
    doi={10.1109/gcat47503.2019.8978281}
    }
    
    @article{ShivaKumarH.2019Lipi,
    title={Lipi Gnani - A Versatile OCR for Documents in any Language Printed in Kannada Script},
    author={R. ShivaKumarH. and A. Ramakrishnan},
    journal={ArXiv},
    year={2019},
    volume={abs/1901.00413},
    doi={10.48550/arxiv.1901.00413}
    }
    
    @article{Singh2026Can,
    title={Can OCR-VLMs Read Devanagari? A Stress-Test Benchmark and Post-Correction Study},
    author={A. Singh},
    journal={ArXiv},
    year={2026},
    volume={abs/2606.29213},
    doi={10.48550/arxiv.2606.29213}
    }
    
    @article{Sivashanth2026Comprehensive,
    title={Comprehensive Evaluation of Tamil OCR Systems: A Survey, Dataset Creation, Benchmarking, and Error Analysis},
    author={S. Sivashanth and Kengatharaiyer Sarveswaran and E. Y. A. Charles},
    journal={International Journal on Advances in ICT for Emerging Regions (ICTer)},
    year={2026},
    doi={10.4038/icter.v18i3.7305}
    }
    
    @article{Taghadouini2026LightOnOCR:,
    title={LightOnOCR: A 1B End-to-End Multilingual Vision-Language Model for State-of-the-Art OCR},
    author={Said Taghadouini and A. Cavaillès and Baptiste Aubertin},
    journal={ArXiv},
    year={2026},
    volume={abs/2601.14251},
    doi={10.48550/arxiv.2601.14251}
    }
    
    @article{Thapa2025Natural,
    title={Natural Language Understanding of Devanagari Script Languages: Language Identification, Hate Speech and its Target Detection},
    author={Surendrabikram Thapa and Kritesh Rauniyar and F. Jafri and Surabhi Adhikari and Kengatharaiyer Sarveswaran and B. Krishna and Bal and Hariram Veeramani and Usman Naseem},
    year={2025},
    pages={71-82}
    }
    
    @article{Tulsyan2024IndicOCR:,
    title={IndicOCR: A Pipeline for Recognizing Printed Documents for Indian Languages},
    author={Krishna Tulsyan and Tessy Flemin and Ajoy Mondal and C. V. Jawahar},
    journal={Proceedings of the 7th Joint International Conference on Data Science & Management of Data (11th ACM IKDD CODS and 29th COMAD)},
    year={2024},
    doi={10.1145/3632410.3632502}
    }
    
    @article{Tulsyan2020A,
    title={A Benchmark System for Indian Language Text Recognition},
    author={Krishna Tulsyan and Nimisha Srivastava and Ajoy Mondal and C. Jawahar},
    year={2020},
    pages={74-88},
    doi={10.1007/978-3-030-57058-3_6}
    }
    
    @article{Vinitha2016Error,
    title={Error Detection in Indic OCRs},
    author={V. Vinitha and C. Jawahar},
    journal={2016 12th IAPR Workshop on Document Analysis Systems (DAS)},
    year={2016},
    pages={180-185},
    doi={10.1109/das.2016.31}
    }
    
    @article{Wang2025Infinity,
    title={Infinity Parser: Layout Aware Reinforcement Learning for Scanned Document Parsing},
    author={Baode Wang and Biao Wu and Weizhen Li and Meng Fang and Yanjie Liang and Zuming Huang and Haozhe Wang and Jun Huang and Ling Chen and Wei Chu and Yuan Qi},
    journal={ArXiv},
    year={2025},
    volume={abs/2510.15349},
    doi={10.48550/arxiv.2506.03197}
    }
    
    @article{Yang2025dots.ocr:,
    title={dots.ocr: Multilingual Document Layout Parsing in a Single Vision-Language Model},
    author={Guang Yang and Hao Liu and Bowen Wang and Colin Zhang},
    journal={ArXiv},
    year={2025},
    volume={abs/2512.02498},
    doi={10.48550/arxiv.2512.02498}
    }
    
    @article{Zilani2026A,
    title={A Hybrid Transformer–Classical Deep Learning OCR Framework for Multilingual Handwritten and Printed Document Recognition},
    author={D. M. Zilani and Mohebbanaaz},
    journal={2026 International Conference on Signal Processing and Electronics Design (ICSPED)},
    year={2026},
    pages={222-228},
    doi={10.1109/icsped65817.2026.11448442}
    }
    
    @article{Zulkarnain2023bbOCR:,
    title={bbOCR: An Open-source Multi-domain OCR Pipeline for Bengali Documents},
    author={Imam Mohammad Zulkarnain and Shayekh Bin Islam and Md. Zami Al Zunaed Farabe and Md. Mehedi Hasan Shawon and Jawaril Munshad Abedin and Beig Rajibul Hasan and M. Haque and Istiak Shihab and Syed Mobassir and Md. Nazmuddoha Ansary and Asif Sushmit and Farig Sadeque},
    journal={ArXiv},
    year={2023},
    volume={abs/2308.10647},
    doi={10.48550/arxiv.2308.10647}
    }
    \end{filecontents}
    
    \documentclass{article}
    \usepackage{adjustbox}
    \begin{document}
    
    \section*{Yes, \textbf{multilingual Indian OCR} now benefits from \textbf{layout-aware parsing, post-correction, and better evaluation}, but same-script language ID and fair benchmarking remain weak spots.}
    
    
    \subsection*{1. Introduction}
    
    
    Multilingual Indian document OCR has moved beyond isolated recognizers toward full pipelines that detect layout, infer reading order, identify script or language, apply OCR, and sometimes reconstruct structured outputs such as searchable HTML, XML, JSON, or field-value records  \cite{Faraz2026Designing} \cite{Tulsyan2024IndicOCR:} \cite{Zulkarnain2023bbOCR:} \cite{Govindaraju2004Tools} \cite{Li2025Printed}. The strongest evidence supports three points: layout analysis materially affects downstream OCR quality, language-model post-correction reduces OCR error for several Indic languages, and evaluation choices such as CER versus WER, Unicode normalization, and real-scan testing can change conclusions about which system is “best”  \cite{Nath2025IndicDLP:} \cite{Bhandari2026A} \cite{Singh2026Can} \cite{Jayatilleke2025Zero-shot}.
    
    
    The evidence is thinner and more mixed for two other issues. Confidence-based fusion of multiple OCR engines often helps in real deployments, but the gains depend on the fusion design and confidence scores are not always reliable proxies for correctness  \cite{Zilani2026A} \cite{Jhala2026PragyaDoc:} \cite{Saluja2017A}. Same-script language identification is now strong for Devanagari-family text in NLP benchmarks, but OCR-specific evidence separating Hindi, Marathi, Sanskrit, and Nepali, or Bengali from Assamese, is much sparser than script-identification evidence across distinct scripts  \cite{Manukonda2025byteSizedLLM@NLU} \cite{Sarungbam2014Script} \cite{Madhani2023Bhasa-Abhijnaanam:} \cite{Mathew2016Multilingual}.
    
    
    \textbf{Figure 1:} Consensus on robust multilingual Indian OCR design
    
    
    \subsection*{2. Methods}
    
    
    This synthesis draws on an expanded Deep Search corpus in Consensus spanning more than 220 million research papers indexed from Semantic Scholar, PubMed, and other sources.
    
    
    The evidence base was screened for papers directly addressing multilingual Indian document OCR pipelines, layout analysis, reading order, OCR engine fusion, post-OCR correction, language or script identification, structured output generation, and evaluation methodology. Priority was given to papers with explicit benchmark results, deployment-oriented system descriptions, or methodological discussion of error metrics and benchmark design.
    
    
    The search combined targeted queries on Indian multilingual OCR, layout and reading-order detection, confidence-based multi-engine fusion, Indic post-OCR correction, same-script language identification, structured OCR outputs, and evaluation validity.
    
    
    \textbf{Figure 2:} Literature search and screening funnel for OCR studies
    
    
    \subsection*{3. Results}
    
    
    \subsubsection*{3.1 Key Papers}
    
    
    A few papers anchor this literature because they connect system design to deployment realities rather than only reporting script-specific recognizer accuracy. Chitrapathak and Parichay frame production-scale multilingual and document-specific OCR for India  \cite{Faraz2026Designing}. IndicDLP establishes that Indic layout parsing lacked adequate multilingual training data until recently  \cite{Nath2025IndicDLP:}. The Devanagari stress-test benchmark is central for evaluation because it shows how clean synthetic benchmarks can misrepresent real-scan performance  \cite{Singh2026Can}.
    
    
    \mbox{}\\
    \begin{adjustbox}{max width=\textwidth}
    \begin{tabular}{p{8.0cm} | p{8.0cm}}
    \hline
    Paper & Summary \\
    \hline
    \cite{Faraz2026Designing} & Production OCR trade-offs for India  \cite{Faraz2026Designing} \\
    \cite{Singh2026Can} & Script-aware Devanagari stress benchmark  \cite{Singh2026Can} \\
    \cite{Nath2025IndicDLP:} & Foundational Indic layout dataset  \cite{Nath2025IndicDLP:} \\
    \hline
    \end{tabular}
    \end{adjustbox}
    
    
    \textbf{Figure 3:} Key papers anchoring multilingual Indian OCR literature
    
    
    \subsubsection*{3.2 Layout and Reading Order}
    
    
    Layout analysis is consistently treated as a prerequisite for accurate OCR in multilingual and mixed-format Indian documents  \cite{Chaudhuri2002Page} \cite{Nath2025IndicDLP:} \cite{Fateh2023Enhancing}. Early Indian systems focused on separating text from non-text, grouping multicolumn text blocks, and handling noisy page structure before recognition  \cite{Chaudhuri2002Page}. Later systems moved toward word- or line-level script handling because whole-line monoscript assumptions fail in real Indian pages  \cite{Pati2004Gabor} \cite{Mathew2016Multilingual} \cite{Kumar2012Discrimination}.
    
    
    \mbox{}\\
    \begin{adjustbox}{max width=\textwidth}
    \begin{tabular}{p{8.0cm} | p{8.0cm}}
    \hline
    Dimension & Evidence \\
    \hline
    Text/non-text separation & Core early requirement  \cite{Chaudhuri2002Page} \cite{Pati2004Gabor} \\
    Multicolumn structure & Explicitly modeled  \cite{Chaudhuri2002Page} \cite{2009} \\
    Reading order & Now supervised explicitly  \cite{Yang2025dots.ocr:} \cite{Li2026Structure-Aware} \\
    Indic data gap & Historically underrepresented  \cite{Nath2025IndicDLP:} \cite{Yang2025dots.ocr:} \\
    OCR impact & Better segmentation raises final accuracy  \cite{Li2025Printed} \cite{Fateh2023Enhancing} \\
    \hline
    \end{tabular}
    \end{adjustbox}
    
    
    \textbf{Figure 4:} Layout and reading-order findings across OCR studies
    
    
    Reading-order detection itself is now treated as a first-class task rather than an afterthought. Unified parsers such as dots.ocr explicitly define document parsing as layout detection, text recognition, and relational understanding through reading-order generation  \cite{Yang2025dots.ocr:}. More general document-parsing work shows that top-down structure constraints or explicit relation prediction improve reading-order robustness on complex layouts  \cite{2009} \cite{Qiao2024Reading} \cite{Li2026Structure-Aware}. Indian OCR papers rarely report a dedicated reading-order metric, but production papers do report failures on dense index pages and complex forms, showing that ordering remains a practical bottleneck  \cite{Kumar2012Discrimination} \cite{Tulsyan2024IndicOCR:}.
    
    
    \subsubsection*{3.3 Fusion, Post-Correction, and Language ID}
    
    
    Confidence-based or consensus-based combination of OCR engines usually beats a weak single engine, but it does not automatically beat the best tuned engine in every setting  \cite{Zilani2026A} \cite{Jhala2026PragyaDoc:} \cite{Saluja2017A}. Hybrid systems using TrOCR, Tesseract, and EasyOCR report higher aggregate multilingual accuracy than any one constituent engine  \cite{Zilani2026A}. In Indian medical documents, confidence-weighted fusion suppressed systematically bad outputs from one engine, but the same study also found that raw OCR confidence was an unreliable surrogate for actual correctness  \cite{Jhala2026PragyaDoc:}. Older Indic correction frameworks likewise found that consolidated outputs from multiple OCR systems reduced required human correction effort  \cite{Saluja2017A}.
    
    
    Language-model post-OCR correction has stronger direct evidence. Context-aware correction reduced CER from 10.20\% to 6.73\% for Hindi, 6.10\% to 1.39\% for Gujarati, and 8.19\% to 3.29\% for Marathi, outperforming seq2seq baselines  \cite{Bhandari2026A}. Sanskrit post-correction benchmarks report a 23-point improvement over OCR output using byte-level tokenization with phonetic encoding  \cite{Maheshwari2022A}. RoundTripOCR extends this to low-resource Devanagari languages by generating millions of synthetic correction pairs across Hindi, Marathi, Nepali, Konkani, Bodo, and Sanskrit, and finds that diverse-font training improves robustness  \cite{Kashid2024RoundTripOCR:}. The main caveat is transfer: a Devanagari byte-level post-corrector improved a cheap engine on matched noise but did not transfer across engines  \cite{Singh2026Can}.
    
    
    Same-script language identification is where OCR pipelines still lean heavily on adjacent NLP work. Native-script IndicLID covers all 22 scheduled languages and is specifically motivated by low performance on similar languages  \cite{Madhani2023Bhasa-Abhijnaanam:}. Devanagari-family identification has become very strong in text benchmarks, with hybrid XLM-R and BiLSTM models reporting F1 0.9974 on Hindi, Marathi, Sanskrit, Nepali, Bhojpuri, and related languages  \cite{Manukonda2025byteSizedLLM@NLU}. For Bengali and Assamese, the literature mostly names the challenge rather than solving it inside OCR pipelines: Bangla script is shared by Bengali, Assamese, and Manipuri, and this similarity complicates language detection beyond script recognition  \cite{Sarungbam2014Script} \cite{Sharma2019Categorization}.
    
    
    \subsubsection*{3.4 Structured Outputs and Evaluation}
    
    
    Structured output support is now common in modern OCR pipelines. IndicOCR exposes script or language ID, segmentation, recognition, and post-correction as pipeline components intended for printed documents in 22 Indian languages  \cite{Tulsyan2024IndicOCR:}. bbOCR reconstructs Bengali pages into structured searchable digital output using layout regions, line and word boxes, and HTML reconstruction  \cite{Zulkarnain2023bbOCR:}. Earlier multilingual annotation platforms already stored Unicode transcriptions in structured XML with metadata and cross-script markup  \cite{Govindaraju2004Tools}. More recent general OCR systems explicitly output editable JSON or Microsoft formats  \cite{Li2025Printed}, while compact OCR-VLMs have started predicting bounding boxes jointly with text  \cite{Taghadouini2026LightOnOCR:} \cite{Hamdi2025PILOT:}.
    
    
    Fair evaluation is the most methodologically mature recent theme. CER and WER measure different failure modes, and low CER can coexist with poor word formation or spacing  \cite{Sivashanth2026Comprehensive} \cite{Jayatilleke2025Zero-shot}. Several recent benchmarks therefore report multiple metrics, including character- and word-level accuracy, normalized edit distance, ANLS, BLEU or chrF++, rather than a single score  \cite{Caravani2026Assessing} \cite{Jayatilleke2025Zero-shot} \cite{Zulkarnain2023bbOCR:}. Unicode handling is not cosmetic in Indic OCR: script-aware evaluation protocols now normalize text to NFC, older robust OCR work identified Unicode re-ordering as a core engineering issue, and synthetic generators increasingly enforce normalization and script purity at data-creation time  \cite{Singh2026Can} \cite{Krishnan2014Towards} \cite{Malik2026synthocr-gen:}.
    
    
    Real-scan validity is another major concern. Clean synthetic text can make systems look nearly tied, while real Devanagari scans cause dramatic separation and expose catastrophic failures, so synthetic-only evaluation overstates quality  \cite{Singh2026Can}. Benchmark builders increasingly try to avoid contamination or circularity by using human-reviewed sources, expert-corrected ground truth, and historically diverse collections  \cite{Faraz2025IndicVisionBench:} \cite{Sivashanth2026Comprehensive}. The strongest benchmarking work also releases tools and configurations for reproducibility and warns that no single OCR paradigm dominates across printed, structured, handwritten, and multilingual settings  \cite{Caravani2026Assessing} \cite{Tulsyan2020A}.
    
    
    \subsubsection*{Results Timeline}
    
    
    Indian multilingual OCR has progressed from script-routing pipelines to layout-aware parsers and OCR-VLMs.
    
    
    \begin{itemize}
    \item \textbf{2002}\begin{itemize}
    \item 1 paper:  \cite{Chaudhuri2002Page}- \textbf{2004}
    \item 2 papers:  \cite{Govindaraju2004Tools} \cite{Pati2004Gabor}- \textbf{2012}
    \item 1 paper:  \cite{Kumar2012Discrimination}- \textbf{2014}
    \item 1 paper:  \cite{Sarungbam2014Script}- \textbf{2016}
    \item 1 paper:  \cite{Mathew2016Multilingual}- \textbf{2017}
    \item 1 paper:  \cite{Saluja2017A}- \textbf{2023}
    \item 3 papers:  \cite{Zulkarnain2023bbOCR:} \cite{Madhani2023Bhasa-Abhijnaanam:} \cite{Fateh2023Enhancing}- \textbf{2024}
    \item 1 paper:  \cite{Tulsyan2024IndicOCR:}- \textbf{2025}
    \item 4 papers:  \cite{Li2025Printed} \cite{Nath2025IndicDLP:} \cite{Jayatilleke2025Zero-shot} \cite{Manukonda2025byteSizedLLM@NLU}- \textbf{2026}
    \item 5 papers:  \cite{Faraz2026Designing} \cite{Bhandari2026A} \cite{Singh2026Can} \cite{Zilani2026A} \cite{Jhala2026PragyaDoc:}\textbf{Figure 5:} Timeline of multilingual OCR progress, larger markers more cited
    \end{itemize}
    
    
    \end{itemize}
    
    
    The arc is clear: early work solved page decomposition, bilingual routing, and Unicode generation; mid-period work emphasized segmentation-free recognizers and error detection; recent work adds layout benchmarks, OCR ensembles, structured outputs, and more rigorous evaluation on real degraded pages  \cite{Chaudhuri1997An} \cite{Krishnan2014Towards} \cite{Vinitha2016Error} \cite{Nath2025IndicDLP:} \cite{Singh2026Can}.
    
    
    \subsubsection*{Top Contributors}
    
    
    Based on visibility within this corpus, \textbf{C. V. Jawahar}, \textbf{A. Ramakrishnan}, and \textbf{U. Pal} appear most frequently across multilingual Indian OCR, script identification, and Indic benchmarking  \cite{Faraz2026Designing} \cite{ShivaKumarH.2019Lipi} \cite{Pal2004Indian}.
    
    
    \mbox{}\\
    \begin{adjustbox}{max width=\textwidth}
    \begin{tabular}{p{5.3cm} | p{5.3cm} | p{5.3cm}}
    \hline
    Type & Name & Papers \\
    \hline
    Author & C. V. Jawahar & [f3cc1bf701865710ba94b827b73ff0bb][f8ff9737fe085c22897d5f12c9d3ecc2][0c0e57467c945c82baa907e7f2673ca8][5d3fd97393185a23b173882da68d2dc9][4796181bda85500f9c4233b300da31b8] \\
    Author & A. Ramakrishnan & [156eb7c0abba576895d763381ec8d5ee][1f0b4e712bba5694baa6bcb401ad1422][d213b76684d75590a829d8a4c3f3c521][ae9218a65fbc5e888b13cf43a171630e] \\
    Author & U. Pal & [0a317776c2b55997b874d54d4ad025d9][330ccbc516095ce4973b483e9c6373bd][d5e4e833ad715f98840cb58037ce538e][88604d05698d55e29d47b4da5880826b] \\
    Journal & \textit{ArXiv} & [f3cc1bf701865710ba94b827b73ff0bb][af7aaaed8c3e516586078a9bb57b398e][1d4ffe7674665b74bbd06206d36fb7cd][14efe715bce05a109b1971b97a56ce7b][7fe0a5f7b529561cb1dbcdf0e1df09f5] \\
    Journal & \textit{IEEE Access} & [722c86ee26365f90a38e946f91f92e0c][7f2ddd0f213a5642bad980b8ce578b01][0b3a505aafdb5f3a8e973866cc8a0137][84bd63e63d2c2d51c5b93bd9da2ebe1a9f] \\
    Journal & \textit{Pattern Recognition} & [88604d05698d55e29d47b4da5880826b][330ccbc516095ce4973b483e9c6373bd][eae59f74267350838abbed6c88ec96e4] \\
    \hline
    \end{tabular}
    \end{adjustbox}
    
    
    \textbf{Figure 6:} Authors and journals appearing most often in corpus
    
    
    \subsection*{4. Discussion}
    
    
    The clearest pattern is that multilingual Indian OCR accuracy now depends less on the recognizer alone and more on pipeline design. Layout analysis, segmentation or parsing, script or language routing, and post-correction all contribute measurable gains, especially on mixed-layout or degraded documents  \cite{Faraz2026Designing} \cite{Li2025Printed} \cite{Tulsyan2024IndicOCR:} \cite{Bhandari2026A}. That is a more mature view than older OCR systems, which largely optimized isolated recognition after hand-designed segmentation  \cite{Sahare2018Multilingual} \cite{Mishra2014WORD} \cite{Ashwin2002A}.
    
    
    The evidence for layout-aware design is strong because it spans old Indian page analyzers, new Indic layout datasets, and current end-to-end parsers that explicitly include reading order as a supervised target  \cite{Chaudhuri2002Page} \cite{Nath2025IndicDLP:} \cite{Wang2025Infinity} \cite{Li2026Structure-Aware}. By contrast, the evidence for same-script language identification inside OCR pipelines is still indirect. The best quantitative results come from NLP shared tasks or native-script LID benchmarks rather than from OCR pages containing mixed Devanagari or Bangla-family content  \cite{Madhani2023Bhasa-Abhijnaanam:} \cite{Manukonda2025byteSizedLLM@NLU} \cite{Thapa2025Natural}.
    
    
    Evaluation practice has improved more rapidly than deployment practice. Recent work increasingly rejects single-metric reporting, adds Unicode normalization, and uses real or historically diverse scans, but many production or demo papers still report word accuracy, exact match, or custom metrics without enough normalization detail to ensure fair comparison  \cite{Singh2026Can} \cite{Caravani2026Assessing} \cite{Tulsyan2020A} \cite{Faraz2026Designing} \cite{Zulkarnain2023bbOCR:}. The Devanagari stress-test is especially important because it shows that benchmark conclusions can flip when moving from rendered pages to real scans and when using median plus catastrophic-failure rate instead of only means  \cite{Singh2026Can}.
    
    
    \mbox{}\\
    \begin{adjustbox}{max width=\textwidth}
    \begin{tabular}{p{4.0cm} | p{4.0cm} | p{4.0cm} | p{4.0cm}}
    \hline
    Claim & Evidence Strength & Reasoning & Papers \\
    \hline
    Layout analysis and reading-order modeling improve multilingual OCR robustness. & Evidence strength: Strong (8/10) & Supported by foundational Indian page-analysis work, new Indic layout resources, and modern parsers with explicit structural supervision. & \cite{Chaudhuri2002Page} \cite{Nath2025IndicDLP:} \cite{Yang2025dots.ocr:} \cite{Li2026Structure-Aware} \\
    Language-model post-OCR correction substantially reduces Indic OCR errors. & Evidence strength: Strong (8/10) & Multiple direct studies report large CER or WER reductions across Hindi, Gujarati, Marathi, Sanskrit, and other Devanagari languages. & \cite{Bhandari2026A} \cite{Maheshwari2022A} \cite{Kashid2024RoundTripOCR:} \cite{Saluja2017Error} \\
    Multi-engine fusion usually beats weak single engines, but confidence alone is insufficient. & Evidence strength: Moderate (6/10) & Several studies show ensemble gains, but one deployment study explicitly finds confidence scores do not reliably indicate correctness. & \cite{Zilani2026A} \cite{Jhala2026PragyaDoc:} \cite{Saluja2017A} \\
    Same-script language ID for Devanagari and Bangla families is solved for OCR pages. & Evidence strength: Weak (3/10) & Strong text-level LID exists, but OCR-page evidence for mixed same-script Indian documents is sparse. & \cite{Manukonda2025byteSizedLLM@NLU} \cite{Madhani2023Bhasa-Abhijnaanam:} \cite{Sarungbam2014Script} \\
    Fair OCR evaluation requires CER/WER, normalization, and real-scan stress tests. & Evidence strength: Strong (9/10) & Recent benchmark papers directly document metric divergence, Unicode issues, and synthetic-to-real performance collapse. & \cite{Sivashanth2026Comprehensive} \cite{Jayatilleke2025Zero-shot} \cite{Singh2026Can} \cite{Caravani2026Assessing} \\
    \hline
    \end{tabular}
    \end{adjustbox}
    
    
    \textbf{Figure 7:} Key claims and supporting evidence across studies
    
    
    \subsection*{5. Conclusion}
    
    
    The literature supports a fairly specific answer: multilingual Indian OCR works best as a layout-aware, script- or language-routed, post-corrected pipeline evaluated on real scans with multiple error metrics, not as a single recognizer judged by word accuracy alone. The most convincing gains are from layout parsing and language-model post-correction, while the weakest evidence concerns same-script language identification inside OCR and the reliability of confidence-only fusion.
    
    
    \subsubsection*{Research Gaps}
    
    
    The strongest work covers printed multilingual OCR, Devanagari post-correction, and broad script identification, but coverage is uneven across same-script language discrimination, reading-order evaluation on Indian pages, and structured-output standards.
    
    
    \mbox{}\\
    \begin{adjustbox}{max width=\textwidth}
    \begin{tabular}{p{2.7cm} | p{2.7cm} | p{2.7cm} | p{2.7cm} | p{2.7cm} | p{2.7cm}}
    \hline
     & Indic Benchmarks & Real Scans & CER Reported & Reading Order & Deployment \\
    \hline
    Layout Parsing & GAP & GAP & GAP & GAP & GAP \\
    Engine Fusion & GAP & GAP & GAP & GAP & GAP \\
    Post-Correction & GAP & GAP & GAP & GAP & GAP \\
    Same-Script LID & GAP & GAP & GAP & GAP & GAP \\
    Structured Output & GAP & GAP & GAP & GAP & GAP \\
    Indic layout resources have improved sharply with IndicDLP, but reading-order evaluation on Indian multilingual pages still lags behind more general document-parsing work  \cite{Nath2025IndicDLP:} \cite{Wang2025Infinity}. Same-script language identification is the starkest gap: Bangla-family similarity is repeatedly recognized as difficult, yet most strong results come from text classification rather than OCR-page routing experiments  \cite{Sarungbam2014Script} \cite{Sharma2019Categorization}. Structured output is increasingly available, but few Indian OCR papers standardize provenance, confidence semantics, or evaluation of JSON or PDF reconstruction quality across systems  \cite{Zulkarnain2023bbOCR:} \cite{Li2025Printed} \cite{Taghadouini2026LightOnOCR:}. \\
    \hline
    \end{tabular}
    \end{adjustbox}
    
    
    \subsubsection*{Open Research Questions}
    
    
    The next advances will likely come from better measurement and better routing rather than only larger recognizers.
    
    
    \mbox{}\\
    \begin{adjustbox}{max width=\textwidth}
    \begin{tabular}{p{8.0cm} | p{8.0cm}}
    \hline
    Question & Why \\
    \hline
    \textbf{Can OCR pipelines reliably distinguish Hindi, Marathi, Sanskrit, Nepali, Bengali, and Assamese directly from noisy page images?} & Same-script LID is central for correct routing, but current evidence is much stronger in clean text benchmarks than in OCR settings. \\
    \textbf{What is the best way to evaluate reading order and structured JSON reconstruction for Indian multilingual documents?} & Layout-aware parsers now output structure, but Indian OCR benchmarks still rarely measure ordering and reconstruction fidelity directly. \\
    \textbf{When does multi-engine fusion outperform the best single OCR model after calibration and post-correction?} & Ensemble gains are real, but confidence scores are imperfect and fair controlled comparisons against tuned single models remain scarce. \\
    \hline
    \end{tabular}
    \end{adjustbox}
    
    
    \textbf{Figure 8:} Open research questions for multilingual Indian OCR
    
    
    The most important shift is from script-specific recognition to document-level parsing and correction; the biggest open question is how to benchmark these richer pipelines fairly on real multilingual Indian pages.
    
    % These search results were found and analyzed using Consensus, an AI-powered search engine for research. Try it at https://consensus.app. © 2026 Consensus NLP, Inc. Personal, non-commercial use only; redistribution requires copyright holders’ consent.
    
    \bibliographystyle{plain}
    \bibliography{refs}
    \end{document}